%% kapa: exam=1 round=36 year=2025 period=1 form=A subject=real_estate subject_ko=부동산학원론 seq=7 number=87
%% source: Q-Net 감정평가사 기출문제 (cst00309) artlSeq=5231525 file=raw/1008/5231525/2025년 제36회 감정평가사 1차시험 1교시 문제지 원본.pdf
%% conversion: kordoc 4.24.51 -> tools/gen1.py (auto); review=proofread-sonnet; flags=none
\begin{question}{87}
부동산 매매시장에서 수요와 공급이 동시에 변화하는 경우, 시장균형의 변화에 관한 설명으로 옳지 않은 것은? (단, 수요곡선은 우하향하고, 공급곡선은 우상향하며, 다른 조건은 동일함)
\begin{choices}
  \item 수요와 공급이 동시에 증가하는 경우, 공급의 증가폭이 수요의 증가폭보다 크면, 균형 가격은 하락하고 균형거래량은 증가한다.
  \item 수요와 공급이 동시에 감소하는 경우, 수요의 감소폭이 공급의 감소폭보다 크면, 균형 가격은 하락하고 균형거래량은 감소한다.
  \item 수요와 공급이 동시에 증가하는 경우, 수요의 증가폭과 공급의 증가폭이 같다면, 균형 가격은 불변이고 균형거래량은 증가한다.
  \item 수요와 공급이 동시에 감소하는 경우, 공급의 감소폭이 수요의 감소폭보다 크면, 균형 가격은 하락하고 균형거래량도 감소한다.
  \item 수요와 공급이 동시에 증가하는 경우, 수요의 증가폭이 공급의 증가폭보다 크면, 균형 가격은 상승하고 균형거래량도 증가한다.
\end{choices}
\end{question}
\answer{4}
